% Worked Examples: Concept 3.9.1 - Modeling Uncertainty
\documentclass[11pt,a4paper]{article}
\usepackage{worked-examples-template}

\title{\textbf{Worked Examples}\\[0.5em]
\large Concept 3.9.1: Modeling Uncertainty\\[0.3em]
\normalsize \coursesubtitle{} -- \coursetitle}
\author{Assoc.\ Prof.\ William Robertson \and Dr Sean McGowan}
\date{\today}

\begin{document}

\maketitle
\tableofcontents
\newpage

\section{Concept 3.9.1: Uncertainty Modelling}

\subsection{Example 1: Parametric Uncertainty in Cruise Control}

Understanding parametric uncertainty is crucial for designing robust control systems. This example demonstrates how parameter variations affect controller performance.

\paragraph{Problem:} 
A cruise control system is designed with a PI controller for a car with nominal mass $m_0 = 1600$ kg. The linearised vehicle dynamics around a constant speed are given by:
\[
m\ddot{x} = F - c\dot{x} - mg\sin\theta
\]
where $x$ is position, $F$ is the drive force, $c = 10$ N·s/m is the drag coefficient, $g = 9.81$ m/s$^2$, and $\theta$ is the road slope.

The PI controller is designed as $C(s) = k_p + k_i/s$ with $k_p = 50$ and $k_i = 20$ for the nominal mass. 

(a) Determine the closed-loop characteristic equation for arbitrary mass $m$.

(b) Calculate the maximum mass $m_{\max}$ for which the closed-loop system remains stable.

(c) Explain why parametric uncertainty matters for robust control design.

\paragraph{Solution:}

\textbf{Step 1: System transfer function}

For velocity control at constant speed, let $v = \dot{x}$ and linearise around a slope angle $\theta$. The transfer function from force to velocity is:
\[
P(s) = \frac{V(s)}{F(s)} = \frac{1}{ms + c}
\]

For $m_0 = 1600$ kg and $c = 10$:
\[
P_0(s) = \frac{1}{1600s + 10} = \frac{1/1600}{s + 0.00625}
\]

\textbf{Step 2: Closed-loop characteristic equation}

With PI controller $C(s) = k_p + k_i/s = (k_ps + k_i)/s$, the loop transfer function is:
\[
L(s) = P(s)C(s) = \frac{1}{ms + c} \cdot \frac{k_ps + k_i}{s} = \frac{k_ps + k_i}{s(ms + c)}
\]

The closed-loop characteristic equation is $1 + L(s) = 0$:
\[
s(ms + c) + k_ps + k_i = 0
\]
\[
ms^2 + (c + k_p)s + k_i = 0
\]

\textbf{Step 3: Stability analysis}

For stability, all coefficients must be positive (necessary condition for 2nd order system):
\begin{itemize}
\item $m > 0$ (always satisfied physically)
\item $c + k_p > 0 \Rightarrow k_p > -c$ (satisfied since $k_p = 50 > 0$)
\item $k_i > 0$ (satisfied since $k_i = 20 > 0$)
\end{itemize}

For a second-order system, if all coefficients are positive, the system is stable. Therefore, the system remains stable for all positive masses.

\textbf{Step 4: Performance degradation}

While the system remains stable, performance degrades with increasing mass. The natural frequency and damping ratio are:
\[
\omega_n = \sqrt{\frac{k_i}{m}}, \quad \zeta = \frac{c + k_p}{2\sqrt{mk_i}}
\]

For $m = 1600$ kg (nominal):
\[
\omega_n = \sqrt{\frac{20}{1600}} = 0.112 \text{ rad/s}, \quad \zeta = \frac{10 + 50}{2\sqrt{1600 \cdot 20}} = \frac{60}{253} = 0.237
\]

For $m = 2000$ kg (25\% increase):
\[
\omega_n = \sqrt{\frac{20}{2000}} = 0.100 \text{ rad/s}, \quad \zeta = \frac{60}{2\sqrt{2000 \cdot 20}} = \frac{60}{283} = 0.212
\]

The response becomes slower (lower $\omega_n$) and slightly less damped (lower $\zeta$), leading to longer settling time and potentially larger overshoot.

\textbf{Key insights:}
\begin{itemize}
\item Parametric uncertainty doesn't necessarily cause instability but degrades performance
\item Controllers designed for nominal conditions may perform poorly under parameter variations
\item Robust design requires considering the full range of expected parameter values
\end{itemize}

\subsection{Example 2: Describing Unmodeled Dynamics with Multiplicative Uncertainty}

Unmodeled dynamics represent system behaviors not captured by our nominal model. This example shows how to characterize them mathematically.

\paragraph{Problem:}
A nominal process model is $P(s) = 1/(s+1)$. However, the actual system has unmodeled high-frequency dynamics and can be represented as $P_{\text{actual}}(s) = 1/[(s+1)(0.1s+1)]$.

(a) Express the actual system using multiplicative uncertainty: $P_{\text{actual}} = P(1 + \delta)$.

(b) Find the multiplicative uncertainty $\delta(s)$ and plot $|\delta(\ii\omega)|$.

(c) Determine the maximum value of $|\delta(\ii\omega)|$ and the frequency at which it occurs.

\paragraph{Solution:}

\textbf{Step 1: Calculate the multiplicative perturbation}

By definition of multiplicative uncertainty:
\[
P_{\text{actual}}(s) = P(s)(1 + \delta(s))
\]

Solving for $\delta(s)$:
\[
\delta(s) = \frac{P_{\text{actual}}(s)}{P(s)} - 1 = \frac{P_{\text{actual}}(s) - P(s)}{P(s)}
\]

Substituting the given models:
\[
\delta(s) = \frac{\frac{1}{(s+1)(0.1s+1)} - \frac{1}{s+1}}{\frac{1}{s+1}} = \frac{1}{(s+1)(0.1s+1)} \cdot (s+1) - 1
\]
\[
= \frac{1}{0.1s+1} - 1 = \frac{1 - (0.1s+1)}{0.1s+1} = \frac{-0.1s}{0.1s+1}
\]

\textbf{Step 2: Frequency response of the uncertainty}

At frequency $\omega$, substitute $s = \ii\omega$:
\[
\delta(\ii\omega) = \frac{-0.1\ii\omega}{0.1\ii\omega + 1} = \frac{-\ii 0.1\omega}{1 + \ii 0.1\omega}
\]

The magnitude is:
\[
|\delta(\ii\omega)| = \left|\frac{-\ii 0.1\omega}{1 + \ii 0.1\omega}\right| = \frac{0.1\omega}{\sqrt{1 + (0.1\omega)^2}}
\]

\textbf{Step 3: Analyze the uncertainty bound}

To find the maximum, take the derivative with respect to $\omega$ and set to zero. Alternatively, observe that:
\[
|\delta(\ii\omega)| = \frac{0.1\omega}{\sqrt{1 + 0.01\omega^2}}
\]

As $\omega \to 0$: $|\delta| \to 0$

As $\omega \to \infty$: $|\delta| \to \frac{0.1\omega}{0.1\omega} = 1$

The function is monotonically increasing and approaches 1 asymptotically.

At $\omega = 10$ rad/s:
\[
|\delta(\ii \cdot 10)| = \frac{1}{\sqrt{1 + 1}} = \frac{1}{\sqrt{2}} = 0.707
\]

At $\omega = 100$ rad/s:
\[
|\delta(\ii \cdot 100)| = \frac{10}{\sqrt{1 + 100}} = \frac{10}{\sqrt{101}} = 0.995
\]

\textbf{Step 4: Physical interpretation}

The multiplicative uncertainty $\delta(s) = -0.1s/(0.1s+1)$ represents:
\begin{itemize}
\item A high-pass characteristic (magnitude increases with frequency)
\item Maximum relative error approaching 100\% at very high frequencies
\item The unmodeled pole at $s = -10$ causes phase lag and gain reduction at high frequencies
\end{itemize}

\textbf{Key insights:}
\begin{itemize}
\item Multiplicative uncertainty naturally increases with frequency for unmodeled fast dynamics
\item The uncertainty bound $|\delta(\ii\omega)| < 1$ for all $\omega$ means the actual plant is within 100\% of the nominal at high frequencies
\item For robust stability, the complementary sensitivity $|T(\ii\omega)|$ must be small where $|\delta(\ii\omega)|$ is large
\end{itemize}
\end{document}
